\documentclass[12pt]{uz-exame} 
\usepackage{amssymb,amsfonts,amsmath}
\newcommand{\rem}[1]{}
\begin{document} 
\examtitle{BSc Honours in Mathematics Level 1}
\examdate{November 2017}
\hours{2}
\lastA{4}
\rubric{\science}
\lastpage{3}
\coursename{CALCULUS 2}
\coursenumber{HMTH111}
%\setcounter{exsec}{17}
\makeheading
\sectionAscience
\begin{question}
Define what is meant by saying that 
\begin{itemize}
\item[(a)] a function $f(x,y)$ has a \textbf{limit}, $L$, at a point $(a,b)$ as $(x,y)$ approaches $(a,b)$. \marks{2}
\item[(b)] a function $f(x,y)$ has a \textbf{partial derivative} $f_x(a,b)$ with respect to $x$ at a point $(a,b)$. \marks{2}
\end{itemize}
\end{question}
\begin{question}
\begin{itemize}
\item[(a)] Find the area of the parallelogram determined by ${\bf{a}}=3{\bf{i}}+4{\bf{j}}$ and ${\bf{b}}={\bf{i}}+{\bf{j}}+{\bf{k}}$.\marks{5}
\item[(b)] Find the volume of a parallelepiped with sides ${\bf a}=3{\bf i}-{\bf j}$, ${\bf b}={\bf j}+2{\bf k}$ and ${\bf c}={\bf i}+5{\bf j}+4{\bf k}$.\marks{6}
\end{itemize}
\end{question}
\begin{question}
Let \[f(x,y) = \left\{ 
\begin{array}{l l}
  \dfrac{xy}{x^2+y^2}, & \quad (x,y)\neq (0,0)\\
  0, & \quad (x,y)=(0,0).\\ \end{array} \right. \]
  \begin{itemize}
  \item[(a)] Determine whether or not $\dfrac{\partial f}{\partial x}$ and $\dfrac{\partial f}{\partial y}$ exist at the origin, and find their values if they exist.\marks{6}
  \item[(b)] Does $\displaystyle \lim _{(x,y)\rightarrow (0,0)} f(x,y)$ exist? \marks{6}
  \end{itemize}
\end{question}
\begin{question}
\begin{itemize}
\item[(a)] Determine whether or not the series converges and find its sum if it converges
\[\sum _{r=1}^\infty \dfrac{1}{(5r-2)(5r+3)}=\dfrac{1}{3\cdot 8}+\dfrac{1}{8\cdot 13}+\dfrac{1}{13\cdot 18}+\cdots +\dfrac{1}{(5n-2)(5n+3)} +\cdots .\]\marks{8}
\item[(b)] By interchanging the order of integration, or otherwise, evaluate the integral \[ \int _{0}^1\int _{3y}^3 e^{x^2}\, dxdy.\]\marks{5}
\end{itemize}
\end{question}
\sectionBscience
\begin{question}
\begin{itemize}
\item[(a)] \begin{itemize}
\item[(i)] Write down the equation of the tangent plane to the paraboloid  $z=x^2+y^2$ at the point $(2,-1,5)$. \marks{5}
\item[(ii)] Find the equation of the plane passing through the points\\ $(2,4,-3), \, (3,7,-1)$ and $(4,3,0)$. \marks{5}
\end{itemize}
\item[(b)]
\begin{itemize}
\item[(i)] Find the angle between the vectors ${\bf{a}}=2{\bf{i}}-3{\bf{j}}+5{\bf{k}}$ and ${\bf{b}}=5{\bf{i}}+3{\bf{j}}-7{\bf{k}}$. \marks{4}
\item[(ii)] Find the parametric equations of the line of intersection of the planes 
\[2x+y+z=4 \quad \mbox{and} \quad 3x-y+z=3\] and calculate the angle between them. \marks{8}
\end{itemize}
\item[(c)] \begin{itemize}
\item[(i)] Find the parametric and symmetric equations for the straight line through $(2,5,-7)$ and $(4,3,8)$. \marks{4}
\item[(ii)] Find the distance between the origin and the plane $x+y+z=10$. \marks{4}
\end{itemize}
\end{itemize}
\end{question}
\begin{question}
\begin{itemize}
\item[(a)] Using the definition of a limit, show that $\displaystyle \lim _{(x,y)\rightarrow (1,2)} x^2+2y=5$. \marks{6}
\item[(b)] Find and classify the stationary points of the function \[f(x,y)=3x-x^3-3xy^2.\] \marks{6}
\item[(c)] Let $w=f(x,y)$, where $x$ and $y$ are given in polar coordinates by the equations $x=r\cos \theta$ and $y=r\sin \theta$.
\begin{itemize}
\item[(i)] Obtain expressions for $\dfrac{\partial w}{\partial r}$ and $ \dfrac{\partial w}{\partial \theta}$.  \marks{8}
\item[(ii)] Show that 
\[\dfrac{\partial ^2 w}{\partial r^2}=\dfrac{\partial ^2 w}{\partial x^2}\cos ^2 \theta +2\dfrac{\partial ^2 w}{\partial x\partial y}\cos \theta \sin \theta +\dfrac{\partial ^2 w}{\partial y^2}\sin ^2 \theta.\]\marks{10}
\end{itemize}
\end{itemize}
\end{question}
\begin{question}
\begin{itemize}
\item[(a)] Sketch the region of integration and evaluate the following integrals.
\begin{itemize}
\item[(i)] $\displaystyle \int _0^1 \int _{x}^{1}e^{-y^2} \, dydx$.
\item[(ii)] $\displaystyle \iint \limits _D \sqrt{x^2+y^2}\, dydx$ where $D$ is the region bounded by $x^2+y^2\geq 4$ and\\ $x^2+y^2\leq 9$. \marks{3,3}
\end{itemize}
\item[(b)] \begin{itemize}
\item[(i)] Evaluate $\displaystyle \int _1^e\int _0^y \int _{0}^{\frac{1}{y}} x^2z\, dxdzdy.$ \marks{5}
\item[(ii)] Evaluate the following integral, using an appropriate transformation or otherwise \[\int _0^1 \int _0^{\sqrt{1-y^2}} \dfrac{dxdy}{1+x^2+y^2}.\] \marks{5}
\end{itemize}
\item[(c)] Let $u=x^2-y^2$ and $v=2xy$.
\begin{itemize}
\item[(i)] Show that $(x^2+y^2)^2=u^2+v^2$. \marks{3}
\item[(ii)] Hence, find $\dfrac{\partial (x,y)}{\partial (u,v)}$ in terms of $u$ and $v$. \marks{5}
\item[(iii)] Hence, evaluate  $\displaystyle \iint \limits _{R} (x^2+y^2)\, dxdy$, where $R$ is the region in the first quadrant  bounded by the hyperbolas $x^2-y^2=6, \, x^2-y^2=1, \, 2xy=4 $ and $2xy=1$. \marks{6}
\end{itemize}
\end{itemize}
\end{question}
\begin{question}
\begin{itemize}
\item[(a)] Test for convergence or divergence for each of the following infinite series by applying an appropriate test.\\
(i)\, $\displaystyle \sum _{n=1}^\infty \dfrac{3^n}{n^2}$ \quad (ii)\, $\displaystyle \sum _{n=1}^\infty \dfrac{\cos n}{n^2}$ \quad (iii)\, $\displaystyle \sum _{n=1}^\infty \dfrac{n}{n^2+1}$. \marks{3,3,3}
\item[(b)] \begin{itemize}
\item[(i)] Find the sum of the series $\displaystyle \sum _{n=1}^\infty\dfrac{1}{n(n+1)}$. \marks{4}
\item[(ii)] Show that the series $\displaystyle \sum _{n=1}^\infty \dfrac{r^n}{n!}$ and $\displaystyle \sum _{n=1}^\infty \dfrac{r^n}{n^n}$ are convergent for all $r$ and that $\displaystyle \sum _{n=1}^\infty n!r^n$ and $\displaystyle \sum _{n=1}^\infty n^n r^n$ are convergent for no $r$ except for $r=0$. \marks{5}
\end{itemize}
\item[(c)] Test the following series for absolute and conditional convergence and determine the interval of convergence.\\
(i)\,  $\displaystyle \sum _{n=0}^\infty \dfrac{2^nx^n}{n!}$ \quad (ii)\, $\displaystyle \sum _{n=1}^\infty \dfrac{x^n}{n3^n} $. 
\marks{6,6}
\end{itemize}
\end{question}
\endexam
\end{document} 